\chapter{Discussion}
\label{ch:discussion}

Chapter~\ref{ch:evaluation} reported the results; this chapter interprets them.
Section~\ref{sec:disc_main} combines the individual results to show where the emotion
route helps and where it does not. Section~\ref{sec:disc_bottleneck} discusses why the
emotion bottleneck helps in those cases, Section~\ref{sec:disc_limitations} states the
limitations of the work, and Section~\ref{sec:disc_relatedwork} relates the findings to
the literature.

\section{Interpreting the Main Result}
\label{sec:disc_main}

The result of this thesis is not a clean win for the emotion route.
Table~\ref{tab:disc_settings} lists the emotion route, direct audio and the PCA-8 control
for every setting in which all three were run, and each row is compared with the chance
floor of its own label space; every score is macro-$F_1$ over films. The emotion route is
clearly ahead of direct audio when the genre classifier is trained on the small Eerola
dataset, both on unseen Eerola films and on the Blockbuster films. It is level with
direct audio when the classifier is trained on the larger Blockbuster dataset, in-domain
there or with both datasets in training, and ahead but not significantly when a
classifier trained on Blockbuster is tested on Eerola.

\begin{table}[t]
  \centering
  \small
  \caption{The three routes in every setting in which all of them were run (film-level macro-$F_1$).}
  \label{tab:disc_settings}
  \setlength{\tabcolsep}{4pt}
  \begin{tabular}{lcrrrrrl}
    \toprule
    Setting & Genres & Floor & Emotion & PCA-8 & Direct & $\Delta$ & $p$ \\
    \midrule
    Eerola, in-domain                        & 5 & 0.301 & 0.469 & 0.418 & 0.375 & $+0.095$ & $0.047$ / $0.151^{\dagger}$ \\
    Blockbuster, in-domain                   & 6 & 0.295 & 0.616 & 0.587 & 0.593 & $+0.023$ & $0.536^{\dagger}$ \\
    Eerola $\rightarrow$ Blockbuster         & 6 & 0.292 & 0.511 & 0.421 & 0.407 & $+0.106$ & $0.018$ \\
    Blockbuster $\rightarrow$ Eerola         & 6 & 0.243 & 0.430 & 0.373 & 0.356 & $+0.077$ & $0.127$ \\
    Both datasets in training                & 6 & --    & 0.529 & 0.539 & 0.553 & $-0.023$ & $0.571^{\dagger}$ \\
    \bottomrule
  \end{tabular}
\end{table}

The rows of Table~\ref{tab:disc_settings} are not comparable with each other, since each
has its own label space, test set and floor (the base-rate random guess). The $p$-value is
the paired film bootstrap unless marked $^{\dagger}$ (Nadeau--Bengio corrected $t$-test over
repeated folds). The Eerola in-domain row gives both; there the Nadeau--Bengio test is
computed on the clip-level scores of the same predictions, since a film-level score per
test fold is not defined. The zero-shot row is the strict (source-scaler) regime, and
Blockbuster rows are films throughout.

\paragraph{In-domain: consistent, larger per film, and partly an operating point.}
On the Eerola dataset the full pipeline (VGGish, predicted emotion, genre) reaches a
film-level macro-$F_1$ of 0.469 on five genres, against 0.385 for the best direct
representation (MusiCNN) and 0.375 for VGGish, its own input. It is ahead of all six
direct representations, by 0.085 to 0.216. This margin is consistent. It holds in all
ten repeats of the cross-validation (nine against the hand-crafted descriptors) and under
all five master seeds of the stability analysis (Section~\ref{sec:eval_robustness}), where
the emotion-minus-VGGish difference ranges from $+0.093$ to $+0.128$ and never changes sign.
The levels are also stable. Across the five seeds, the standard deviation of each route's
ten-repeat mean is at most 0.015, so the reported levels reflect the dataset rather than one
random split. On eight genres the advantage over both direct representations run there
is larger still ($+0.096$ over VGGish and $+0.116$ over AST) and significant at film level and,
over clips, under both tests used in this thesis.

However, on five genres its significance depends on the comparison and the test. The film
bootstrap estimates how much the result depends on which films make up the dataset. At
film level it finds the advantage significant against VGGish ($p=0.047$) and against AST,
CLAP, the hand-crafted descriptors and wav2vec~2.0, but not against MusiCNN
($p=0.079$). The Nadeau--Bengio test also accounts for the variation that comes from which
films the models were trained on. It cannot be applied per film, and on the clip-level
scores of the same predictions it finds the advantage over VGGish borderline
($p=0.151$). The two tests do not contradict each other, because they measure
different sources of uncertainty. Together they indicate that the in-domain advantage is
real in its direction, but that its size is not yet established.

The metric analysis of Section~\ref{sec:eval_robustness} qualifies this result further.
The emotion route ranks first under macro-$F_1$ over films and over clips, weighted $F_1$,
macro average precision and macro ROC-AUC, so its lead depends neither on one averaging
scheme nor on the unit of scoring. If each genre's decision threshold is instead tuned on
the training fold, all routes converge over clips (0.427--0.441) and the advantage disappears ($+0.011$, $p=0.439$ against
VGGish; $-0.001$, $p=0.924$ against the PCA-8 control). Under the two threshold-free
ranking metrics it is small and not significant. Part of the in-domain advantage at the
default threshold of one half is therefore an \emph{operating-point} effect. With
class-balanced weighting, the classifier on eleven emotion features places more of the
true positives above the default cut and so reaches a much higher recall at only slightly
higher precision. The direct classifiers rank clips almost as well, but cut them in the
wrong place. This result supports two readings. In the narrow reading, the emotion route
does not rank genres better than direct audio in-domain; its predicted probabilities are
better placed relative to the default cut. In practical terms, a model that works at
the default threshold without per-genre tuning has a real advantage, since such tuning
needs labelled data from the same distribution and becomes unreliable for genres
represented by a handful of films. This advantage is, however, smaller and more specific
than ``emotion is better''. Film-level scoring uses the same default cut on averaged
probabilities and therefore inherits this caveat; thresholds were not tuned per film.

\paragraph{Across datasets: an advantage that depends on the training data.}
When the model is trained on Eerola and applied without any further fitting to the 110
Blockbuster films, the emotion route reaches 0.511 against 0.407 for direct audio
($+0.106$, $p=0.018$), and 0.511 against 0.421 for the PCA-8 control ($+0.092$, $p=0.032$).
The reverse direction, trained on Blockbuster and tested on the 41 Eerola films, agrees in
sign ($+0.077$) but is not significant ($p=0.127$). When both datasets are included in the
training data, the gap disappears ($-0.023$, $p=0.571$). In-domain on Blockbuster alone, the
emotion route is also level with direct audio (Table~\ref{tab:disc_settings}).

The pattern across the five rows of Table~\ref{tab:disc_settings} is the main result of
this thesis. The emotion route is clearly ahead where the genre classifier has been
trained on the 41 Eerola films, whether it is tested on unseen Eerola films or on the
Blockbuster films, and level where it has been trained on the 110 Blockbuster films or on
both datasets. The pooled design shows both sides at once: on the Eerola films of its
test folds the emotion route leads (0.537 against 0.413), and on the Blockbuster films
direct audio does (0.573 against 0.539). Its advantage is therefore an advantage
\emph{under scarce data and under a change of dataset}. An eight-dimensional
representation whose axes mean the same thing in every dataset needs few examples and
survives a change of dataset. A 128-dimensional embedding, whose axes carry the statistics
of the data it was computed on, needs more examples, and catches up once the classifier
is trained on enough of them or on the target distribution itself. This matches what is
expected of an interpretable intermediate, and it is a more defensible claim than
``emotion is better at predicting genre''.

The zero-shot comparison is a single train--test split and answers a different question
from the in-domain threshold analysis. In a zero-shot setting no labelled target data is
available, so no threshold can be tuned on it, and the default operating point is the one a
practitioner would use. The operating-point pattern is nevertheless visible here as well.
From Eerola to Blockbuster the emotion route gains its lead through much higher recall at
lower precision, which is a stronger form of the recall-driven pattern seen in-domain.
Whether the zero-shot advantage survives thresholds tuned on the \emph{source} dataset was
not tested and would still be useful to know (Section~\ref{sec:futurework}). However, an
operating-point effect is unlikely to explain the whole advantage. In the reverse
direction, scored per film, the emotion route is ahead of direct audio in both precision
(0.375 against 0.360) and recall (0.622 against 0.481), although the difference in $F_1$ is
not significant there. A shift
of the decision threshold alone cannot produce this, because it trades one for the other.

\section{Why the Bottleneck Helps When It Helps}
\label{sec:disc_bottleneck}

\subsection{Meaning rather than compression}

A sceptical reading of the result is that eleven numbers simply overfit less than 128 on a
dataset of about 330 clips from 41 films. Three pieces of evidence speak against this
reading, although the first only in part. The PCA-8 control is just as compact. Over
clips it is no better than the embedding it compresses in-domain ($-0.005$, $p=0.771$),
and under zero-shot transfer it gains only $+0.014$ ($p=0.572$). At film level, however,
it is ahead of VGGish in-domain under every seed ($+0.033$ to $+0.059$), so compression does help
there, but the emotion route stays ahead of the control under every seed as well
($+0.040$ to $+0.070$). Compression alone therefore does not explain the advantage. Nor is
the control a weak baseline. On Blockbuster it
beats a random eight-dimensional projection of the same embedding by a wide margin, so it
keeps the most informative directions that a generic method can find. Second, the ablation
shows that the genre-relevant content of the emotions is low-dimensional and redundant
(Section~\ref{subsec:disc_ablation}). The emotion bottleneck could therefore be narrower
still, which makes an eight-dimensional control a conservative comparison. Finally, the
emotion--genre signatures derived from human ratings on Eerola reappear on Blockbuster,
where they are computed from the emotions that the Eerola-trained model predicts for each
cue ($r=0.836$ over 48 genre--emotion cells). In both datasets Horror is defined by fear.
A principal component has no meaning that could replicate in this way, and an artefact of
the Eerola dataset would not reappear in 110 other films.

The evidence against the control is, however, weaker than the evidence against direct
audio. At the default threshold the emotion route is ahead of the PCA-8 control in every
design except the pooled one, where the two are level, but significantly in only two
cases: in the strict zero-shot experiment ($p=0.032$) and in-domain over clips under the
film bootstrap ($p=0.029$, Nadeau--Bengio $p=0.122$). At film level the in-domain lead is
$+0.051$ ($p=0.326$). When the transfer was repeated in both directions with one shared
implementation (Section~\ref{sec:eval_transfer}), the emotion route was again ahead of the
control, but not significantly ($+0.045$, $p=0.258$ from Eerola to Blockbuster;
$+0.058$, $p=0.145$ in reverse). This was largely because the control scored higher there than in the original
zero-shot run. Single-split scores of the embedding arms moved by up to 0.04 between two
legitimate implementations, while the emotion arm remained stable. With per-genre
thresholds the two bottlenecks tie in-domain. The claim that the gain comes from the
meaning of the eight numbers rather than from their number is therefore supported in
direction in every design except the pooled one. It is significant in the strict
zero-shot run and, in-domain, only over clips under the film bootstrap, and it is
supported most strongly by the replication of the signatures. It is not supported with the same force as the claim against direct audio.

\subsection{Why more information does not win}

The direct embeddings contain strictly more information about the audio than the eight
emotions predicted from them, since every emotion prediction is a function of the
embedding. There are three explanations, which reinforce each other, for why they
nevertheless do worse on the small Eerola dataset and across datasets, and catch up only
when the classifier is trained on the larger Blockbuster dataset.

The first is \emph{data scarcity}. With about 330 clips from about 40 films, a linear
classifier over 128 to 768 dimensions can fit many directions that happen to separate the
genres of the training films. Strong regularisation was selected for almost every feature
set, which is a symptom of this problem, and the learning curves of both routes were still
rising at the full dataset size. The emotion stage adds prior knowledge that the direct
classifier would otherwise have to learn from very few examples. It was trained on
listener ratings, which are a second and independent source of supervision, and it reduces
the embedding to the directions that listeners consider emotionally relevant. This
explanation fits the pattern of Table~\ref{tab:disc_settings}, in which the direct
embedding catches up exactly where more training films are available.

A second explanation is \emph{film identity}. Clips from one film share their
instrumentation, recording and mixing, and a rich embedding encodes these properties well.
Under a naive split this makes the direct baseline look far stronger than it is, because
the model recognises the film instead of the genre. Under the film-grouped protocol used
throughout this thesis, the same information is at best useless and at worst misleading,
since it rewards learning the sound of individual training films. The emotion features are
much less sensitive to this effect, because an emotion rating describes what a clip
expresses rather than how it was produced.

The third explanation is \emph{domain shift}. An embedding's dimensions have no fixed
meaning, and their scale and distribution depend on the recordings they were computed on.
A classifier trained on them partly learns the idiosyncrasies of its training dataset,
which is harmless in-domain but harmful when the dataset changes. Emotion values are
expressed in the units of a rating scale and mean the same thing in any dataset, so the
representation absorbs the change of dataset. The pooled design shows this most clearly.
Once both datasets are in the training data, the direct classifier can learn both sets of
idiosyncrasies, and its disadvantage disappears.

\subsection{A low-dimensional genre signal}
\label{subsec:disc_ablation}

The ablation study returned a null result, but one that is still informative. No single
emotion is necessary. Removing any one of them changes the film-level score by at most
0.010, and none of these changes is significant. Smaller subsets lose a little: fear alone,
or valence and energy alone, keep 93 to 95\,\% of the score of all eight, lower in
nearly every repeat but not significantly. The reason lies in the ratings themselves. Fear
rises with tension and falls with valence, and the genre signal they carry lies along
essentially one threat--versus--warmth axis, which a small subset of emotions can largely
recover.

For the genre task, this means that the part of film music that distinguishes the five
genres modelled here is low-dimensional. This is one reason why richer representations
cannot extract more from it with the data available. As for the design of the thesis, all
eight emotions are retained because the discrete categories make the model
readable (Section~\ref{subsec:fund_why_not_va}), even though the classifier does not need
them. ``Horror is fear'' is a better description of a score than ``Horror is low valence
and high tension'', although a classifier can use either.

\subsection{The Stage-1 ceiling}

On paper, the emotion regressor is the weakest part of the pipeline. The best
representations explain about 0.56 of the variance in the listener ratings, which is about
62\,\% of the variance that two independent listener panels agree on
(Section~\ref{sec:eval_emotion}). Even so, genre predicted from these imperfect emotions comes
close to genre predicted from the human ratings. At film level the ratings are ahead by
0.053 (0.469 against 0.521), in every repeat but just short of significance ($p=0.053$), and over
clips the two are level (0.417 against 0.428, $p=0.325$). The ablation explains why the
gap is so much smaller than the gap in $R^2$. The genre signal lies in the threat--warmth
axis of fear, tension and valence, which are among the better-predicted emotions. Happiness and sadness are the hardest to predict (the
categorical form of the ``valence problem'' of the emotion-recognition literature), but
given fear, tension and valence they are largely redundant for the classifier. A better
emotion model would therefore improve genre prediction, as the film-level gap shows, but
by less than its $R^2$ gap suggests, because the genre-relevant part of Stage~1 is closer to
its ceiling than the average.

The waveform representation provides a counterweight.
wav2vec~2.0, the weakest representation used directly, gains most from the bottleneck
($+0.222$ at film level and $+0.091$ over clips, significant under every test) and reaches the
same level as the VGGish pipeline.
Part of the effect of the bottleneck is therefore \emph{denoising}. Forcing a poor
representation through a supervised, low-dimensional target discards the directions that
carry nothing about genre. This effect is real but not specific to interpretation, and it
explains why the gain is largest for the representation with the least to lose.

\subsection{What the bottleneck buys in readability}

Apart from accuracy, the emotion route provides what it was designed for. Each genre the
model can predict has a readable affective signature: Horror is fear, Comedy and Drama are
high valence and low threat, and Action is tension without happiness. Adventure, the one
genre without a signature, is also predicted poorly for its frequency (Biography and
Documentary fail because there are too few films), and the pipeline makes
the reason for this failure visible. This is the practical value of an interpretable
intermediate. Both its errors and its successes come with an explanation that can be
checked against film-music practice.

\section{Limitations}
\label{sec:disc_limitations}

\paragraph{Dataset size.} The Eerola dataset contains 360 rated clips from 46 films, 346 of them (43 films) with a modelled genre (41
films in the five-genre subset), and this size limits every in-domain conclusion. The
learning curves are still rising at the full dataset size. For every borderline
five-genre comparison, the smallest $p$-value that the Nadeau--Bengio test could reach on
the clip-level scores with infinitely many repeats is above 0.05 ($p_{\lim}$ between 0.055
and 0.130), so additional computation cannot settle these comparisons; only more films
can. Rare genres make the problem worse. Comedy is represented by four films and is
absent from a third of the individual test folds. Since the headline score is computed
over films, the $F_1$ of Comedy and Horror rests on four and five films (Documentary and
Biography on two and three), which is why the film-level scores spread two to three times
as much as the clip-level ones.

\paragraph{Two tests that disagree.} At film level only the film bootstrap can be
computed. Over clips, on five genres, the film bootstrap finds the in-domain advantage
significant and the Nadeau--Bengio test finds it borderline. Both are reported because
they measure different things. The first measures the dependence on which films
make up the dataset, and the second also the dependence on which films a model was trained
on. Reporting only the more favourable of the two would overstate the result. Because the
zero-shot experiment is a single split, it can only be tested with the film bootstrap.

\paragraph{Dependence on the operating point.} The in-domain advantage holds at the default
decision threshold, where the emotion route ranks first under macro-$F_1$ and four further
metrics. However, its lead is significant under only some of them, it is small and not
significant under the threshold-free ones, and it disappears with per-genre thresholds
tuned on the training data, where all routes converge. The headline in-domain comparison
is therefore a statement about models used at a fixed, untuned operating point. In
addition, several multi-label metrics reward always predicting the frequent genres
(micro-$F_1$, sample-averaged $F_1$, exact match, Hamming loss and Jaccard), and the
most-frequent baseline wins them outright. For this reason macro-$F_1$ is the headline
metric, and these metrics cannot be used to rank the routes at all.

\paragraph{Choice of the five-genre subset.} The five headline genres were selected partly
because they carry an emotional signature in the same data, which can favour the emotion
route. The eight-genre results, in which the advantage over AST and VGGish is significant
at film level and, over clips, under both tests, reduce this risk but do not remove it.

\paragraph{Unit of evaluation.} Scoring per film was adopted after the clip-level analysis
had been carried out, when the analysis of alternative metrics showed that it fits the
structure of the labels. A metric chosen after the results are known can favour the
method it is chosen for, so the clip-level score of the same predictions is reported
alongside throughout. The main conclusions hold under both. The exceptions are the
comparisons with MusiCNN and with the PCA-8 control, which are significant over clips
under the film bootstrap but not over films; the comparison of predicted with rated
emotions, which is level over clips but favours the ratings over films; and the reverse
transfer direction, which is significant only under a bootstrap that resamples clips.

\paragraph{Genre labels.} Genre labels come from IMDb and are a property of the film, not
of the individual clip or cue. A romantic scene in an action film is labelled Action, and a
comic interlude in a drama is labelled Drama, so some clips carry labels that do not
describe the music they contain. Genre is also multi-label, and IMDb lists a film's genres
alphabetically rather than by relevance, so there is no ``primary'' genre that could be
used to simplify the task or to weight the labels. The resulting label noise is probably
the largest single source of error, and it caps what any audio model can achieve.

\paragraph{Perceived emotion, one population.} The emotion model learns \emph{perceived}
emotion, that is, what listeners judged the music to express rather than the emotion it
induced in them. It learns this from the ratings of one listener panel. A second panel
that rated a subset of the clips agreed closely with the first, but both were recruited
for the same study. Their agreement therefore bounds the noise among comparable listeners
and says nothing about listeners from a different musical culture. The film-music codes
described by \citet{gorbman1987unheard} are learned conventions, and such listeners might
rate the same clips differently.

\paragraph{The second dataset.} The Blockbuster dataset distributes neither audio (for
copyright reasons) nor emotion ratings. As a result, only VGGish, the embedding published
with the dataset, could be carried across, so the transfer result rests on a single
representation. In addition, the emotions on Blockbuster are predicted by the
Eerola-trained model and cannot be checked directly against human judgement. They are
validated only indirectly, through the replication of the emotion--genre signatures and
through the transfer result itself. Finally, the zero-shot comparison is a single split
whose point estimates move by a few hundredths with legitimate implementation choices, so
the result rests on its confidence intervals rather than on its third decimal.

\section{Relation to Prior Work}
\label{sec:disc_relatedwork}

The comparison with \citet{ma2021computational} is like-for-like and does not rely on an
unverified baseline. Their protocol, reproduced on their own published features, reaches
0.620 in-domain, inside their reported range of 0.61--0.65. The reproduction also qualifies
their observation that VGGish generally outperforms the hand-crafted features. Under
repeated cross-validation the learned embedding does not beat their full set of
hand-crafted features ($+0.009$, $p=0.787$), and the two are statistically
indistinguishable. This agrees with Ma et al.'s own cautious wording and with the result
on Eerola, where direct representations of very different design score close to one
another. Across two datasets, the choice of audio
representation therefore matters less than what is done with it.

\citet{eerola2011genrespecific} found that models of musical emotion lose much of their
accuracy when carried from one musical genre to another. This thesis observes the same
degradation one level higher, where genre models lose accuracy when carried from one film
dataset to another, as the Related Work chapter anticipated. The present results add that
the loss is smaller for a model built on emotions than for one built on raw embeddings.
This suggests that the repertoire dependence described by
\citeauthor{eerola2011genrespecific} lies more in the mapping from audio to its
representation than in the relationship between emotion and genre itself.

The other genre studies are consistent with this. \citet{austin2010characterization} found
high-intensity genres more separable by their scores than narrative genres. The signatures
found here explain this, since Horror carries the strongest emotional signature, Action a distinct one, and
Adventure has none. \citet{mangolin2020multimodal} found audio to be the weakest of five
modalities for genre, and the modest absolute scores reported here agree that audio alone
is a limited predictor of genre. The contribution of this thesis lies in a model whose
predictions can be explained and which degrades less when the data change, rather than in
a higher number. This also answers the concern raised by \citet{kang2025there} that
accurate emotion-recognition systems frequently reveal little about the mechanisms behind
their predictions. When predicted emotion is used as an intermediate for a downstream task,
the emotion model becomes an explanation instead of an end in itself. Finally, the metric
analysis supports the decision of this thesis not to follow \citet{behrouzi2023multimodal}
in using Hamming loss as a headline metric, because on these data a classifier that always
predicts the most frequent genres wins on Hamming loss outright.
